% answers_short.tex —— 2016 年·上海卷·理科 · 简短答案速查
% 由 tools/make_short_answers.py 自 answers.tex 的 \gksolhead/\gkanswer 提取，
% 并与 2016/2016_数学真题_上海理_答案要点.md 逐题比对后生成，勿手改（改源 MD 后重跑）。

\begin{gknotice}
\textbf{使用说明}：本节只给最终答案便于快速核对，完整推导见解析版。答案为编写者独立推导，\textbf{非官方标准答案}。
\end{gknotice}

\gkpart{填空题}
\begin{itemize}[leftmargin=2.6em]
  \item[\textbf{1.}] \((2,4)\)
  \item[\textbf{2.}] \(-3\)
  \item[\textbf{3.}] \(\dfrac{2}{\sqrt5}\)
  \item[\textbf{4.}] \(1.76\) 米
  \item[\textbf{5.}] \(f^{-1}(x)=\log_2(x-1)\)，\(x>1\)
  \item[\textbf{6.}] \(2\sqrt2\)
  \item[\textbf{7.}] \(\dfrac{\pi}{6},\dfrac{5\pi}{6}\)
  \item[\textbf{8.}] \(112\)
  \item[\textbf{9.}] \(\dfrac{7\sqrt3}{3}\)
  \item[\textbf{10.}] \((2,+\infty)\)
  \item[\textbf{11.}] \(4\)
  \item[\textbf{12.}] \([0,1+\sqrt2]\)
  \item[\textbf{13.}] \(4\)
  \item[\textbf{14.}] \(\dfrac{5}{28}\)
\end{itemize}

\gkpart{选择题}
\begin{center}\begin{tabular}{ccccc}
\toprule
题号 & 15 & 16 & 17 & 18 \\
\midrule
答案 & A & D & B & D \\
\bottomrule
\end{tabular}\end{center}

\gkpart{解答题}
\begin{itemize}[leftmargin=2.6em]
  \item[\textbf{19.}] (1) \(\dfrac{\sqrt3}{12}\)；(2) \(\dfrac{\pi}{4}\)
  \item[\textbf{20.}] (1) \(y^2=4x\)，\(0\le y\le2\)；(2) 矩形面积 \(\dfrac52\)，五边形面积 \(\dfrac{11}{4}\)，五边形更接近 \(\dfrac83\)
  \item[\textbf{21.}] (1) \(y=\pm\sqrt2\,x\)；(2) \(k=\pm\sqrt{\dfrac35}\)
  \item[\textbf{22.}] (1) \((-\infty,-\dfrac14)\cup(0,+\infty)\)；(2) \((1,2]\cup\{3,4\}\)；(3) \([\dfrac23,+\infty)\)
  \item[\textbf{23.}] (1) \(a_3=16\)；(2) 不具有性质 P；(3) 对任意初值都具有性质 \(P\) 的充要条件是 \(\{b_n\}\) 为常数列
\end{itemize}
